\begin{afterword}
\summaryhead

The author recalls being a young ``Traditional Rationalist'' who gave a mysterious answer to a mysterious question, the most embarrassing mistake of the author's life. The answer is not described, but it made a bold prediction: that neurologists would find neurons exploiting quantum gravity, after Roger Penrose. The young author obeyed the traditional rules. The theory was falsifiable, avoided magic and called itself physics. Yet it explained nothing, and after it the phenomenon was as mysterious as before.

A Bayesian, the author says, would have asked why the theory deserved belief over ``I don't know,'' since it made no retrospective predictions. Traditional Rationality helped the author out of the error but could not get things right. It is taught through physicists' biographies, would accept thirty years spent on a silly idea, and lets people guess and agree to disagree. The Way of Bayes is also imprecise, but it has mathematics and experimental psychology behind it.

\responsehead

The post is candid about its author's error, and it grants two things a reader might otherwise hold against it: that Traditional Rationality helped the author climb out of the error, and that the Bayesian alternative is also an imprecise art. The trouble is with the case it argues from and with the contrast it draws.

The case is withheld. The post will not describe the ``mysterious answer'' because ``it would be a long tale and complicated.'' It gives only the prediction, which was Penrose's, and never names the question, which was consciousness. So its diagnosis, complicated thinking without every step pinned down, cannot be checked against the theory. From this one episode the post draws a general law: you need ``one whole hell of a lot of rationality'' before it does more than lead you into new mistakes. No other case is offered.

The two traditions are drawn to fit the story. ``Traditional Rationality'' is the author's own category, described from biographies of physicists; no text or teacher is quoted. It is said to ignore retrospective predictions, but science has counted them: the old data on Mercury's orbit confirmed general relativity, and it is Bayesian theories of confirmation that have had trouble saying why. The tradition is said to lack the ideal of ``only one correct probability estimate given the evidence,'' but that ideal divides Bayesians too, since the subjective school permits any coherent prior. And it is said that the tradition would have let the author spend thirty years on the idea. In practice, the physical idea the author borrowed was argued over by calculation, without waiting for an experiment.

The post does not show the Bayesian method at work. It names the question a Bayesian would ask, whether the hypothesis beats ``I don't know,'' but does not ask it of the theory, which the reader is not shown. Its support for the new method is ``underlying math'' and experimental psychology, named but not cited.

In short: a candid confession used as evidence for a general law, with the case withheld and a contrast between two traditions drawn more sharply than the history of science, or of Bayesianism, allows.
\end{afterword}
